\documentclass[../../main.tex]{subfiles}

\begin{document}

\section{紧拓扑空间}\label{section:ch11-紧拓扑空间}

我们已对度量空间研究了紧性. 我们给出了紧性的几个刻画且证明了连续映射与定义在紧度量空间上的连续函数的性质. 紧性的概念可被自然且有用地推广到拓扑空间.

回忆集族 $\{E_{\lambda}\}_{\lambda\in\Lambda}$ 称为是集合 $E$ 的覆盖, 若 $E\subseteq\bigcup_{\lambda\in\Lambda}E_{\lambda}$. 若每个 $E_{\lambda}$ 包含于一个拓扑空间, 一个覆盖称为是开的, 若该覆盖中的每个集合是开的.

\begin{definition}[紧空间]\label{definition:ch11-紧空间}
拓扑空间 $X$ 称为紧的, 若 $X$ 的每个开覆盖有有限子覆盖. $X$ 的子集 $K$ 称为紧的, 若 $K$ 作为具有继承自 $X$ 的子空间拓扑的拓扑空间是紧的.
\end{definition}

根据子空间拓扑的定义, $X$ 的子集 $K$ 是紧的, 若 $K$ 的每个由 $X$ 的开子集族组成的覆盖有有限子覆盖.

关于拓扑空间的紧性的某些结果可直接从度量空间移植过来, 例如, 紧空间在连续映射下的象也是紧的. 紧度量空间的其他性质, 例如, 紧性与序列紧性的等价性仅可移植到具有某些有附加拓扑结构的拓扑空间. 紧度量空间的另外一些性质, 如全有界性, 在拓扑背景下没有简单的对应结果.

回忆集族称为具有有限交性质, 若它的每个有限子族有非空交. 由于拓扑空间 $X$ 的子集是闭的当且仅当它的补在 $X$ 中是开的, 根据 De Morgan 等式, 我们先前对度量空间证明过的一个结果在拓扑空间有如下推广.

\begin{proposition}\label{proposition:ch11-有限交}
拓扑空间 $X$ 是紧的当且仅当 $X$ 的每个具有有限交性质的闭子集族有非空交.
\end{proposition}

\begin{proposition}\label{proposition:ch11-紧子集闭}
紧拓扑空间 $X$ 的闭子集 $K$ 是紧的.
\end{proposition}

\begin{proof}
令 $\{\mathcal{O}_{\lambda}\}_{\lambda\in\Lambda}$ 为 $K$ 的由 $X$ 中的开子集组成的开覆盖. 由于 $X\sim K$ 是 $X$ 的开子集, $[X\sim K]\cup\{\mathcal{O}_{\lambda}\}_{\lambda\in\Lambda}$ 是 $X$ 的开覆盖. 根据 $X$ 的紧性该覆盖有有限子覆盖, 通过可能地从该有限子覆盖中去掉集合 $X\sim K$, 剩下的族是覆盖 $K$ 的 $\{\mathcal{O}_{\lambda}\}_{\lambda\in\Lambda}$ 的有限子族. 因此 $K$ 是紧的.
\end{proof}

\begin{proposition}\label{proposition:ch11-Hausdorff紧闭}
Hausdorff 拓扑空间 $X$ 的紧子空间 $K$ 是 $X$ 的闭子集.
\end{proposition}

\begin{proof}
我们将证明 $X\sim K$ 是开的, 因此 $K$ 必须是闭的. 令 $y$ 属于 $X\sim K$. 由于 $X$ 是 Hausdorff 的, 对每个 $x\in K$ 分别存在 $x$ 和 $y$ 的不交邻域 $\mathcal{O}_{x}$ 和 $U_{x}$. 则 $\{\mathcal{O}_{x}\}_{x\in K}$ 是 $K$ 的开覆盖. 由于 $K$ 是紧的, 存在有限子覆盖 $\{\mathcal{O}_{x_{1}},\mathcal{O}_{x_{2}},\cdots,\mathcal{O}_{x_{n}}\}$. 定义 $N=\bigcap_{i=1}^{n}U_{x_{i}}$. 则 $N$ 是 $y$ 的邻域, 它与每个 $\mathcal{O}_{x_{i}}$ 不交, 因此包含于 $X\sim K$. 因此 $X\sim K$ 是开的.
\end{proof}

\begin{definition}[序列紧]\label{definition:ch11-序列紧}
拓扑空间 $X$ 称为序列紧的, 若 $X$ 中的每个序列有一个收敛到 $X$ 中的点的子序列.
\end{definition}

我们已证明了度量空间是紧的当且仅当它是序列紧的. 对第二可数拓扑空间同样的结果成立.

\begin{proposition}\label{proposition:ch11-第二可数紧}
令 $X$ 为第二可数拓扑空间. 则 $X$ 是紧的当且仅当它是序列紧的.
\end{proposition}

\begin{proof}
首先假设 $X$ 是紧的. 令 $\{x_{n}\}$ 为 $X$ 中的序列. 对每个指标 $n$, 令 $F_{n}$ 为非空集 $\{x_{k}\mid k\geqslant n\}$ 的闭包. 则 $\{F_{n}\}$ 是一个递减的非空闭集序列. 由于 $\{F_{n}\}$ 具有有限交性质, 根据命题 14, $\bigcap_{n=1}^{\infty}F_{n}\neq\varnothing$, 选取这个交集中的点 $x_{0}$. 由于 $X$ 是第二可数的, 它是第一可数的. 令 $\{B_{n}\}_{n=1}^{\infty}$ 为拓扑在点 $x_{0}$ 的基. 我们可以假设每个 $B_{n+1}\subseteq B_{n}$. 由于对每个 $n$, $x_{0}$ 属于 $\{x_{k}\mid k\geqslant n\}$ 的闭包, $B_{n}$ 的邻域与 $\{x_{k}\mid k\geqslant n\}$ 有非空交. 因此可以归纳地选取严格递增的指标序列 $\{n_{k}\}$ 使得对每个指标 $k$, $x_{n_{k}}\in B_{k}$. 由于对 $x_{0}$ 的每个邻域 $\mathcal{O}$, 存在指标 $N$ 使得对 $n\geqslant N$, $B_{n}\subseteq\mathcal{O}$, 子序列 $\{x_{n_{k}}\}$ 收敛到 $x_{0}$. 因此 $X$ 是序列紧的.

现在假定 $X$ 是序列紧的. 由于 $X$ 是第二可数的, 每个开覆盖有可数子覆盖. 因此为证明 $X$ 是紧的仅须证明 $X$ 的每个可数覆盖有有限子覆盖. 令 $\{\mathcal{O}_{n}\}_{n=1}^{\infty}$ 为这样的一个覆盖. 我们用反证法. 假设没有有限子覆盖. 则对每个指标 $n$, 存在指标 $m(n)>n$, 使得 $\mathcal{O}_{m(n)}\sim\bigcup_{i=1}^{n}\mathcal{O}_{i}\neq\varnothing$. 对每个自然数 $n$, 选取 $x_{n}\in\mathcal{O}_{m(n)}\sim\bigcup_{i=1}^{n}\mathcal{O}_{i}$. 则由于 $X$ 是序列紧的, $\{x_{n}\}$ 的一个子序列收敛到 $x_{0}\in X$. 但 $\{\mathcal{O}_{n}\}_{n=1}^{\infty}$ 是 $X$ 的一个开覆盖, 因此存在某个 $\mathcal{O}_{N}$ 是 $x_{0}$ 的邻域. 因此, 有无穷多个指标 $n$ 使得 $x_{n}$ 属于 $\mathcal{O}_{N}$. 这是不可能的, 由于对 $n>N$, $x_{n}\notin\mathcal{O}_{N}$.
\end{proof}

\begin{theorem}\label{theorem:ch11-紧Hausdorff正规}
紧 Hausdorff 空间是正规的.
\end{theorem}

\begin{proof}
令 $X$ 为紧与 Hausdorff 的. 我们首先证明它是正则的, 即每个闭集合与不在该集合中的点可被不交的邻域分开. 令 $F$ 为 $X$ 的闭子集而 $x$ 属于 $X\sim F$. 由于 $X$ 是 Hausdorff 的, 对每个 $y\in F$ 存在 $x$ 与 $y$ 的不交邻域 $\mathcal{O}_{x}$ 与 $\mathcal{U}_{y}$. 则 $\{\mathcal{U}_{y}\}_{y\in F}$ 是 $F$ 的开覆盖. 但 $F$ 是紧的. 因此有有限子覆盖 $\{\mathcal{U}_{y_{1}},\mathcal{U}_{y_{2}},\cdots,\mathcal{U}_{y_{n}}\}$. 定义 $N=\bigcap_{i=1}^{n}\mathcal{O}_{y_{i}}$. 则 $N$ 是 $y$ 的邻域, 它与 $F$ 的邻域 $\bigcup_{i=1}^{n}\mathcal{U}_{y_{i}}$ 不交. 因此 $X$ 是正则的. 重复这个方法可证明 $X$ 是正规的.
\end{proof}

\begin{proposition}\label{proposition:ch11-紧到Hausdorff一一同胚}
紧空间 $X$ 映上 Hausdorff 空间 $Y$ 的连续一一映射 $f$ 是一个同胚.
\end{proposition}

\begin{proof}
为证明 $f$ 是同胚, 仅须证明它将开集映到开集或等价地将闭集映到闭集. 令 $F$ 为 $X$ 的闭子集. 则 $F$ 是紧的, 由于 $X$ 是紧的. 因此, 根据命题 20, $f(F)$ 是紧的. 因此, 根据命题 16, 由于 $Y$ 是 Hausdorff, $f(F)$ 是闭的.
\end{proof}

\begin{proposition}\label{proposition:ch11-紧连续象}
紧拓扑空间的连续映射下的象是紧的.
\end{proposition}

\begin{proof}
令 $f$ 为紧拓扑空间 $X$ 到拓扑空间 $Y$ 的连续映射. 令 $\{\mathcal{O}_{\lambda}\}_{\lambda\in\Lambda}$ 为 $f(X)$ 的由 $Y$ 的开子集组成的覆盖. 则根据 $f$ 的连续性, $\{f^{-1}(\mathcal{O}_{\lambda})\}_{\lambda\in\Lambda}$ 是 $X$ 的一个开覆盖. 根据 $X$ 的紧性, 存在覆盖 $X$ 的有限族 $\{f^{-1}(\mathcal{O}_{\lambda_{1}}),\cdots,f^{-1}(\mathcal{O}_{\lambda_{n}})\}$. 有限族 $\{\mathcal{O}_{\lambda_{1}},\cdots,\mathcal{O}_{\lambda_{n}}\}$ 覆盖 $f(X)$.
\end{proof}

\begin{proposition}\label{proposition:ch11-紧到Hausdorff}
映上 Hausdorff 空间 $Y$ 的紧空间 $X$ 的连续一一映射 $f$ 是同胚映射.
\end{proposition}

\begin{proof}
为了证明 $f$ 是同胚映射, 仅须证明它将开集带入开集, 或者等价地它将闭集带入闭集. 令 $F$ 是 $X$ 的闭集. 则 $F$ 是紧的, 因为 $X$ 是紧的. 因此, 根据命题 20, $f(F)$ 是紧的. 于是, 根据命题 16, 由于 $Y$ 是 Hausdorff 空间, $f(F)$ 是闭的.
\end{proof}

\begin{corollary}\label{corollary:ch11-紧空间最值}
紧拓扑空间上的连续实值函数取到最大值与最小值.
\end{corollary}

\begin{proof}
令 $X$ 为紧的而 $f:X\to\mathbf{R}$ 是连续的. 根据前一个命题, $f(X)$ 是实数的一个紧集. 因此 $f(X)$ 是闭与有界的. 而实数的闭有界集包含最小与最大成员.
\end{proof}

拓扑空间称为可数紧的, 若每个可数开覆盖有一个有限子覆盖. 我们在习题 39 与 40 探索这样的空间的性质.

\begin{exercise}\label{exercise:ch11-39}
对于第二可数空间 $X$, 证明 $X$ 是紧的当且仅当它是可数紧的.
\end{exercise}

\begin{exercise}\label{exercise:ch11-40}
(Fr\'echet 交定理)令 $X$ 为拓扑空间. 证明 $X$ 是可数紧的当且仅当 $\{F_{n}\}$ 是递减的 $X$ 的非空闭子集序列, 交 $\bigcap_{n=1}^{\infty}F_{n}$ 是非空的.
\end{exercise}

\begin{exercise}\label{exercise:ch11-41}
令 $X$ 为紧的、Hausdorff 的, 而 $\{F_{n}\}_{n=1}^{\infty}$ 是 $X$ 的下降的闭子集族. 令 $\mathcal{O}$ 为交 $\bigcap_{n=1}^{\infty}F_{n}$ 的邻域. 证明存在指标 $N$ 使得对 $n\geqslant N$, $F_{n}\subseteq\mathcal{O}$.
\end{exercise}

\begin{exercise}\label{exercise:ch11-42}
证明不可能把实数的闭有界区间表示为闭有界区间的可数族(多于一个成员)的两两不交并.
\end{exercise}

\begin{exercise}\label{exercise:ch11-43}
令 $f$ 为紧空间 $X$ 映上 Hausdorff 空间 $Y$ 的连续映射. 证明任何 $Y$ 映入 $Z$ 的使得 $g\circ f$ 连续的映射 $g$ 自身是连续的.
\end{exercise}

\begin{exercise}\label{exercise:ch11-44}
令 $(X,\mathcal{T})$ 为拓扑空间.
\begin{enumerate}[(i)]
\item 证明: 若 $(X,\mathcal{T})$ 是紧的, 则对任何比 $\mathcal{T}$ 弱的拓扑 $\mathcal{T}_{1}$, $(X,\mathcal{T}_{1})$ 是紧的.
\item 证明: 若 $(X,\mathcal{T})$ 是 Hausdorff 的, 则对任何比 $\mathcal{T}$ 强的拓扑 $\mathcal{T}_{2}$, $(X,\mathcal{T}_{2})$ 是 Hausdorff 的.
\item 证明: 若 $(X,\mathcal{T})$ 是紧与 Hausdorff 的, 则任何严格弱的拓扑不是 Hausdorff 的, 而任何严格强的拓扑不是紧的.
\end{enumerate}
\end{exercise}

\begin{exercise}\label{exercise:ch11-45}
(紧开拓扑)令 $X$ 与 $Y$ 为 Hausdorff 拓扑空间, 而 $Y^{X}$ 是从 $X$ 到 $Y$ 的映射族. 在 $Y^{X}$ 上, 通过取形如 $\mathcal{U}_{K,\mathcal{O}}=\{f:X\to Y\mid f(K)\subseteq\mathcal{O}\}$ 的集合为子基, 我们定义了称为紧开拓扑的拓扑, 其中 $K$ 是 $X$ 的紧子集而 $\mathcal{O}$ 是 $Y$ 的开子集. 因此紧开拓扑是 $Y^{X}$ 上使得集合 $\mathcal{U}_{K,\mathcal{O}}$ 是开的最弱拓扑.
\begin{enumerate}[(i)]
\item 令 $\{f_{n}\}$ 为 $Y^{X}$ 中的序列, 关于紧开拓扑收敛到 $f\in Y^{X}$. 证明在 $X$ 上 $\{f_{n}\}$ 逐点收敛到 $f$.
\item 现在假设 $Y$ 是度量空间. 证明 $Y^{X}$ 中的序列 $\{f_{n}\}$ 关于紧开拓扑收敛到 $f\in Y^{X}$, 当且仅当在 $X$ 的每个紧子集 $K$ 上 $\{f_{n}\}$ 一致收敛到 $f$.
\end{enumerate}
\end{exercise}

\begin{exercise}\label{exercise:ch11-46}
(Dini 定理)令 $\{f_{n}\}$ 为可数紧空间 $X$ 上的连续实值函数序列. 假定对每个 $x\in X$, 序列 $\{f_{n}(x)\}$ 单调递减到零. 证明 $\{f_{n}\}$ 一致地收敛到零.
\end{exercise}

\end{document}